\section*{Bab \ref{ch:b3:continuum-elasticity} --- Mekanika Kontinum dan Elastisitas}

\begin{solution}{exo:b3:continuum-elasticity:1}
(a) $\unit{Pa} = \unit{N/m^2} = \unit{kg.m^{-1}.s^{-2}}$; regangan
adalah bilangan murni. (b) $\varepsilon = \sigma/E = \num{e-3}$. (c)
Melipatduakan panjangnya berarti $\varepsilon = 1$: seribu kali di luar
``kecil'', dan tanggapan karet di sana sangat taklinear (dan berasal
dari sebab lain, yakni entropi) --- \omterm{thm:b3:continuum-elasticity:hooke}{hukum Hooke} hanya membuka kurvanya.
(d) $w =
\sigma^2/2E$: baja $(\num{2e8})^2/2(\num{2e11}) =
\qty{1e5}{J/m^3}$; karet $(\num{e6})^2/2(\num{2e6}) =
\qty{2.5e5}{J/m^3}$ --- bahan yang lunak menyimpan \emph{lebih banyak}
per satuan volume pada tegangan kerjanya, dan itulah sebabnya ketapel
dibuat dari karet, bukan dari baja.
\end{solution}

\begin{solution}{exo:b3:continuum-elasticity:2}
(a) $\varepsilon_{ij} = \alpha\delta_{ij}$: \omterm{def:b3:continuum-elasticity:strain}{dilatasi} isotropik,
$\delta V/V = 3\alpha$. (b) $\varepsilon_{xy} = \varepsilon_{yx} =
\gamma/2$, sisanya nol: geseran murni, volumenya tak berubah. (c)
$\varepsilon_{ij} = 0$: rotasi tegar, tanpa deformasi. (d)
$\operatorname{diag}(\varepsilon, -\nu\varepsilon, -\nu\varepsilon)$:
uji tarik --- pemuluran sepanjang $x$, penciutan Poisson melintangnya.
\end{solution}

\begin{solution}{exo:b3:continuum-elasticity:3}
(a) Dengan hanya $\sigma_{zz}(z)$ dan beratnya: $\partial_z\sigma_{zz} =
\rho g$ (dengan mampatan diambil positif ke bawah), $\sigma_{zz} = \rho
gh$ di dasarnya. (b) $h_{\max} = \sigma_{\text{c}}/\rho g =
\num{2e8}/(2700 \times 9.81) \approx \qty{7.6}{km}$. (c) Everest berada
tepat pada batas hancur dasarnya sendiri --- gunung di Bumi setinggi
yang diizinkan kekuatan batuan, dan tak lebih tinggi. (d) Kerucut
setinggi $h$ membebani dasarnya hanya dengan $\rho gh/3$ (sepertiga
tiangnya: massanya tumbuh bersama penampangnya), sehingga timbunan
berbentuk gunung dapat berdiri kira-kira tiga kali lebih tinggi
daripada sebuah tiang.
\end{solution}

\begin{solution}{exo:b3:continuum-elasticity:4}
(a) $-\nu\varepsilon = \num{-3e-4}$. (b) $\delta V/V = (1 -
2\nu)\varepsilon = \num{4e-4}$. (c) $\nu = 1/2$: deformasi tak
termampatkan --- karet. (d) Untuk $\nu > 1/2$ modulus ruahnya $K =
E/3(1 - 2\nu)$ akan negatif: menekan dari segala sisi justru akan
\emph{menumbuhkan} volumenya, dan bahannya akan melepaskan energi dengan
meruntuh --- tak ada zat padat mantap yang dapat melakukannya.
\end{solution}

\begin{solution}{exo:b3:continuum-elasticity:5}
(a) Pada separuh silinder sepanjang satu satuan, tekanannya mendorong
dengan $p
\times 2R$ (luas proyeksinya) dan kedua dinding irisnya menarik balik
dengan $2\sigma_\theta t$: $\sigma_\theta = pR/t$. (b) Pada sebuah
penampang, $p\pi R^2 = \sigma_z\,2\pi Rt$: $\sigma_z = pR/2t$ ---
separuhnya. Sosis pecah \emph{memanjang}: tegangan lingkarnya yang lebih
besar merobek kulitnya sepanjang sumbu. (c) $\sigma_\theta = \num{2e7} \times
0.09/0.005 = \qty{360}{MPa}$: sebesar faktor $2$ di bawah luluhnya ---
dan itulah sebabnya tabung diuji tekan dan diperiksa berkala. (d) Bola
memikul $pR/2t$ ke segala arah: ujung setengah bola menyeparuhkan
tegangannya dan menghindari sudut, tempat ujung datar akan memusatkannya.
\end{solution}

\begin{solution}{exo:b3:continuum-elasticity:6}
(a) $G = E/2(1 + \nu) = \qty{1.0}{MPa}$. (b) $\sigma_{xy} = F/S =
150/\num{e-2} = \qty{15}{kPa}$; $\gamma = \sigma_{xy}/G = 0.015$;
perpindahannya $\gamma h = \qty{0.3}{mm}$. (c) Jaringan karet mengubah
\emph{bentuk} hanya dengan mengatur ulang arah rantainya (mudah),
tetapi mengubah \emph{volume} berarti memampatkan rantainya lebih rapat,
sesukar memampatkan zat cair: $G \ll K$, sehingga $\nu \to 1/2$. (d)
Kelunakan terhadap geseran: tumpuannya membiarkan tanah berguncang
mendatar di bawahnya sambil meneruskan sedikit saja gaya, namun tetap
cukup kaku terhadap mampatan untuk memikul berat gedungnya.
\end{solution}

\begin{solution}{exo:b3:continuum-elasticity:7}
(a) $k = ES/L = \num{1.2e9} \times \num{8e-5}/20 = \qty{4.8}{kN/m}$.
(b) $mg(h + x) = \tfrac12 kx^2$: $2400x^2 - 785x - 3140 = 0$, $x =
\qty{1.3}{m}$. (c) $F = kx \approx \qty{6.3}{kN}$; perlambatannya $F/m
\approx \qty{79}{m/s^2} \approx 8g$ --- itulah sebabnya tali panjat
sengaja dibuat mulur. (d) Sebagian energi yang terserap masuk ke
pemutusan serat dan penataan ulang plastis: talinya kini berada pada
kurva tegangan--regangan yang merosot, lebih kaku dan lebih lemah, dan
jatuh berikutnya akan lebih keras dalam kedua arti itu.
\end{solution}

\begin{solution}{exo:b3:continuum-elasticity:8}
(a) Puncak tabung berjejari $r$ berputar sebesar busur $r\theta$
sepanjang $L$: $\gamma = r\theta/L$. (b) $\Gamma = \int_0^a r\,(G r\theta/L)\,
2\pi r\,\dd r = (\pi Ga^4/2L)\,\theta$. (c) Sebesar $2^4 = 16$. (d) $C =
\pi \times \num{4e10} \times (\num{2.5e-5})^4/2 =
\qty{2.5e-8}{N.m/rad}$; $T = 2\pi\sqrt{I/C} = 2\pi\sqrt{\num{5e-4}/
\num{2.5e-8}} \approx \qty{900}{s}$ --- seperempat jam per ayunan:
kepekaan yang cukup untuk merasakan gravitasi bola timbal.
\end{solution}

\begin{solution}{exo:b3:continuum-elasticity:9}
(a) $\mu = E/2(1+\nu) = \qty{78}{GPa}$, $\lambda = E\nu/(1+\nu)(1-2\nu)
= \qty{107}{GPa}$: $c_{\text{P}} = \sqrt{262\times10^9/7850} =
\qty{5.8}{km/s}$, $c_{\text{S}} = \qty{3.1}{km/s}$. (b) Batang tipis
memberikan $\sqrt{E/\rho} = \qty{5.0}{km/s}$: pada batangnya sisi-sisinya
bebas menggembung (pelepasan Poisson), yang melunakkan tanggapannya;
dalam ruahnya bahan di sekitarnya melarangnya. (c) $c_{\text{S}} = 0$;
$c_{\text{P}} = \sqrt{\num{2.2e9}/1000} = \qty{1.5}{km/s}$ --- yakni
laju bunyi dalam air. (d) P menggerakkan tanah sepanjang sinarnya
(hampir tegak untuk sumber dalam), S melintanginya: ketiga komponen
seismometer memisahkan polarisasinya, dan selang P/S lalu menetapkan
jaraknya.
\end{solution}

\begin{solution}{exo:b3:continuum-elasticity:10}
(a) Busur pada jejari $R + y$ berpanjang $(R + y)\alpha$ berbanding
$R\alpha$ pada permukaan netralnya: $\varepsilon = y/R$, $\sigma =
Ey/R$. (b) $M = \int y\sigma\,\dd S = (E/R)\int y^2\,\dd S = EI/R$;
untuk persegi panjang $I = bh^3/12$. (c) Rebah: $I = bh^3/12$; tegak:
$hb^3/12$ --- nisbahnya $(b/h)^2 = 25$ untuk $b/h = 5$: kayu yang sama,
dua puluh lima kali lebih kaku, semata-mata karena geometrinya. (d) $I = 0.50 \times
0.035^3/12 = \qty{1.8e-6}{m^4}$, $EI = \qty{2.1e4}{N.m^2}$; $\delta =
FL^3/3EI = 736 \times 1.8^3/(3 \times \num{2.1e4}) \approx
\qty{7}{cm}$: persis kelenturan papan loncat yang nyaman itu.
\end{solution}

\begin{solution}{exo:b3:continuum-elasticity:11}
(a) Dalam keadaan melengkung, beban aksial $P$ bekerja dengan lengan
gaya $y(x)$ terhadap penampang di $x$: $M = -Py$, dan $M = EIy''$
memberikan $EIy'' = -Py$. (b) $y = A\sin(x\sqrt{P/EI})$ dengan
$y(L) = 0$: $A$ tak nol pertama kali muncul pada
$\sqrt{P/EI}\,L = \pi$, yakni $P_{\text{c}} =
\pi^2EI/L^2$; di bawahnya tiang lurus adalah satu-satunya penyelesaian,
di atasnya melengkung tidak menuntut gaya sama sekali. (c) $I = 0.03 \times 0.0015^3/12 =
\qty{8.4e-12}{m^4}$: $P_{\text{c}} = \pi^2 \times 12\times10^9 \times
\num{8.4e-12}/1^2 \approx \qty{1}{N}$ --- seberat sebuah apel: penggaris
satu meter tertekuk oleh satu jari. (d) $I$ tumbuh ketika bahannya
bergerak menjauhi sumbunya: sebuah tabung menempatkan seluruh
penampangnya pada $y$ besar, sehingga $I$ (dan karenanya
$P_{\text{c}}$) maksimum untuk berat bahan tertentu --- itulah rancangan
tulang, bambu dan kerangka sepeda.
\end{solution}

\begin{solution}{exo:b3:continuum-elasticity:12}
(a) Satu rantai per luas $a^2$; meregangkannya sebesar regangan
$\varepsilon$ meregangkan tiap pegasnya sebesar $\delta = \varepsilon a$,
dengan gaya $\kappa
\varepsilon a$ per rantai dan tegangan $\kappa\varepsilon/a$: $E =
\kappa/a$. (b) $\rho = m/a^3$, sehingga $E/\rho = \kappa a^2/m$ dan $c =
a\sqrt{\kappa/m}$. (c) $c = \num{2.5e-10}\sqrt{50/\num{e-25}} \approx
\qty{5.6}{km/s}$. (d) Tepat pada beberapa km/s yang terukur untuk logam:
kekakuan segala sesuatu yang padat --- dan laju setiap gelombang gempa
--- adalah kekakuan ikatan kimianya.
\end{solution}

\begin{solution}{pb:b3:continuum-elasticity:1}
\textbf{1.} $\mu = E/2(1+\nu) = \qty{30}{GPa}$; $\lambda =
E\nu/(1+\nu)(1-2\nu) = \qty{30}{GPa}$: untuk $\nu = 1/4$, $\lambda =
\mu$.
\textbf{2.} $c_{\text{P}} = \sqrt{3\mu/\rho} = \sqrt{\num{9e10}/2700}
= \qty{5.8}{km/s}$; $c_{\text{S}} = \sqrt{\mu/\rho} =
\qty{3.3}{km/s}$; nisbahnya $\sqrt3$.
\textbf{3.} $\varepsilon \sim 2/\num{5e4} = \num{4e-5}$; $\sigma =
E\varepsilon \approx \qty{3}{MPa}$ --- ``penurunan tegangan'' yang
terukur pada gempa nyata memang beberapa MPa.
\textbf{4.} $w = \tfrac12 E\varepsilon^2 \approx \qty{60}{J/m^3}$;
pada $\num{5e13}{}\,\unit{m^3}$: $\sim\qty{3e15}{J}$, sekitar
tiga perempat megaton --- gempa yang besar.
\textbf{5.} Dari bawah, P menggerakkan tanah secara tegak --- gedung
dibangun untuk memikul beban tegak; S menggerakkannya secara
\emph{mendatar}, yakni arah gedung paling lemah.
\textbf{6.} Hanya \omterm{prop:b3:continuum-elasticity:waves}{gelombang P}: ia gelombang mampatan; \omterm{prop:b3:continuum-elasticity:waves}{gelombang S}
adalah geseran murni, pada volume tetap.
\textbf{7.} $t_{\text{P}} = d/c_{\text{P}}$, $t_{\text{S}} =
d/c_{\text{S}}$: $\Delta t = d(1/c_{\text{S}} - 1/c_{\text{P}})$,
yang dibalik seperti dinyatakan.
\textbf{8.} $1/(1/3.33 - 1/5.77)\,\unit{km/s} \approx
\qty{7.9}{km}$ per sekon selangnya --- ``delapan kilometer per sekon''
yang dipakai ahli gempa di lapangan.
\textbf{9.} $d \approx 7.9 \times 28 \approx \qty{220}{km}$.
\textbf{10.} Tiga (pada umumnya): tiap stasiun mengenal satu lingkaran
episentrum yang mungkin; dua lingkaran berpotongan di dua titik, dan
yang ketiga memutuskan.
\textbf{11.} $t_{\text{P}} = 80/5.8 = \qty{14}{s}$, $t_{\text{S}} =
80/3.3 = \qty{24}{s}$: sekitar \qty{10}{s} peringatan antara sentakannya
dan guncangan yang merusak.
\textbf{12.} S tiba pada $200/3.3 = \qty{60}{s}$. P mencapai
\qty{20}{km} pada \qty{3.5}{s}; pesan radio praktis seketika, sehingga
kotanya dapat memperoleh hampir satu menit peringatan --- sistem
sekarang mencapai puluhan sekon.
\textbf{13.} $2R_{\text{E}}/c \approx 12742/10 \approx \qty{1270}{s}$,
sekitar dua puluh menit.
\textbf{14.} Bahwa \omterm{prop:b3:continuum-elasticity:waves}{gelombang elastik} melintasinya sama sekali: Bumi
dalam tidak lebur seluruhnya --- ia meneruskan gelombang, dan (untuk
\omterm{prop:b3:continuum-elasticity:waves}{gelombang S} yang melintasi mantel) bahkan tergeser, seperti zat padat.
\textbf{15.} Dua jejari dan sudut $\Delta$ di antaranya: tali busurnya
adalah $2R_{\text{E}}\sin(\Delta/2)$.
\textbf{16.} $2 \times 6371 \times \sin30^\circ = \qty{6371}{km}$;
$t \approx \qty{640}{s} \approx \qty{11}{min}$.
\textbf{17.} Di luar $103^\circ$ setiap sinar harus melewati sebuah
daerah pusat yang dalam; jika daerah itu \emph{fluida} ($\mu = 0$), ia
tidak memikul gelombang geser: \omterm{prop:b3:continuum-elasticity:waves}{gelombang S} mati di sana. Fluida adalah
wujud zat yang tidak dapat melawan geseran.
\textbf{18.} Fluida tetap memikul mampatan: \omterm{prop:b3:continuum-elasticity:waves}{gelombang P} melintasi
intinya, tetapi lebih lambat --- sehingga membias tajam di batasnya, dan
sinar yang membelok itu meninggalkan bayangan berbentuk cincin, bukan
potongan yang bersih.
\textbf{19.} Sebuah inti cair berada di pusat Bumi.
\textbf{20.} Jarak terdekat tali busurnya ke pusat adalah
$R_{\text{E}}\cos(\Delta/2)$; menyerempet berarti
$R_{\text{E}}\cos(\Delta^*/2) = R_{\text{c}}$, yakni $\Delta^* =
2\arccos(R_{\text{c}}/R_{\text{E}})$.
\textbf{21.} $R_{\text{c}} = 6371\cos(51.5^\circ) \approx
\qty{3970}{km}$.
\textbf{22.} Terlalu besar $+14\%$. Lajunya tumbuh seiring kedalaman,
sehingga sinar yang sebenarnya melengkung kembali ke arah permukaan:
sinar yang titik terdalamnya baru menyentuh intinya muncul pada sudut
yang \emph{lebih kecil} daripada tali busur lurus yang menembus
kedalaman yang sama --- jadi dari $103^\circ$ yang teramati, pembalikan
sinar lurus menempatkan titik singgungnya terlalu dangkal, sehingga
menaksir intinya terlalu besar.
\textbf{23.} Energi P yang lemah di dalam bayangannya berarti ada
sesuatu di dalam inti cair yang membelokkan sinarnya keluar kembali:
yakni sebuah \emph{inti dalam} tersendiri
(Lehmann, 1936).
\textbf{24.} Gelombang geser hanya ada dalam zat padat: \omterm{prop:b3:continuum-elasticity:waves}{gelombang P}
yang berubah menjadi gelombang geser di dalam inti dalamnya lalu
kembali (fase yang disebut PKJKP) akan membuktikan bahwa ia padat ---
deteksinya, yang lama dicari, adalah bukti yang diterima.
\textbf{25.} Dua laju dalam batuan ($5.8$ dan \qty{3.3}{km/s}), satu
gelombang yang hilang di luar $103^\circ$, dan geometri tali busur
menghasilkan inti cair sebesar $\approx\qty{4000}{km}$ --- dalam
selisih $15\%$ dari nilai modern
\qty{3480}{km}: \omterm{thm:b3:continuum-elasticity:hooke}{hukum Hooke}, yang dibaca pada skala planet.
\end{solution}
